% Standalone source-record documentation; no research attempts.
\documentclass[11pt]{article}
\usepackage[margin=25mm]{geometry}
\usepackage{fontspec,unicode-math}
\setmainfont{Latin Modern Roman}
\setmathfont{Latin Modern Math}
\usepackage{amsmath,amssymb,mathtools,xurl,hyperref}
\hypersetup{colorlinks=true,urlcolor=blue,linkcolor=blue}
\setcounter{secnumdepth}{0}
\setlength{\parindent}{0pt}
\setlength{\parskip}{5pt}
\setlength{\emergencystretch}{3em}
\newcommand{\sref}[2]{\hyperlink{source:#1:#2}{[S#2]}}
\newcommand{\cataloguescope}[1]{\paragraph{Source record.}#1}
\newcommand{\floor}[1]{\left\lfloor#1\right\rfloor}
\begin{document}
% UnsolvedMath ID 30003378; source catalogue code OWR-15210-004
\setcounter{section}{30003377}
\section{Error-Free Multiparty Homomorphic Secret Sharing}\label{30003378}
{\small\sffamily UnsolvedMath ID 30003378 · Computer Science}
\subsection{Definitions and mathematical statement}

\paragraph{Shared notation.}Standard number systems are $\mathbb N,\mathbb Z,\mathbb Q,\mathbb R,\mathbb C$. Specialized notation and index ranges follow the cited source and the definitions below.


Homomorphic secret sharing distributes an input among parties so that local computations yield shares of the target output. A branching program is a directed decision graph evaluating a function by following input-labelled transitions. DDH is the decisional Diffie-Hellman assumption in the selected group.
\paragraph{Record scope and qualifications.}
The catalogue title is retained as metadata. Named certificates, contracts, bounds or frameworks in that title are not asserted as results of the original source extract.
The error convention, compactness target, and group/security models are not specified in the short question.

Source-record limitation: The sharing security model, branching-program class, compactness condition, and error/running-time requirement are not specified.
\paragraph{Statement recorded in the supplied source.}
For DDH-based homomorphic secret sharing supporting compact evaluation of branching programs, can the inverse-polynomial evaluation error be eliminated or its running-time dependence improved, and can analogous results be obtained for three or more parties? \sref{30003378}{2}
\paragraph{Source status.}
The supplied catalogue labels this source record partially solved; that metadata was checked against the supplied snapshot on 2026-10-11. The label is preserved for provenance and is not an independent verification of present open status or mathematical correctness. The catalogue literature assessment is dated 2026-08-21.
\subsection{Short English statement}
Remove or improve the stated homomorphic-sharing error and extend the result to three or more parties.
\subsection{Sources}
\begin{itemize}
\item[\hypertarget{source:30003378:1}{S1}] Ulam AI and the UnsolvedMath Contributors, UnsolvedMath: A Curated Collection of Open Mathematics Problems, record 30003378 (OWR-15210-004). Catalogue curation and metadata: CC BY 4.0. The source extract is retained; definitions, English summaries and qualifications are editorial additions.
\url{https://huggingface.co/datasets/ulamai/UnsolvedMath}.
\item[\hypertarget{source:30003378:2}{S2}] Cryptography (2018), Oberwolfach Reports, DOI 10.4171/owr/2017/5, pp. 213--266
\url{https://doi.org/10.4171/owr/2017/5}.
\item[\hypertarget{source:30003378:3}{S3}] Abram et al., Succinct Homomorphic Secret Sharing
\url{https://eprint.iacr.org/2024/814.pdf}.
\item[\hypertarget{source:30003378:4}{S4}] Multi-Party Homomorphic Secret Sharing
\url{https://eprint.iacr.org/2023/1593.pdf}.
\item[\hypertarget{source:30003378:5}{S5}] Homomorphic Secret Sharing
\url{https://eprint.iacr.org/2018/419.pdf}.
\item Source-specific Oberwolfach Reports licences, including CC BY-SA 4.0 where applicable, are retained; see THIRD\_PARTY\_NOTICES.md.
\end{itemize}
\paragraph{Preparation and review.}
The original source and UnsolvedMath contributors retain attribution for the source extract. Codex (GPT-6; exact serving revision not exposed) assisted with definitions, English summaries, scope notes and formatting. No new human mathematical authorship or independently verified solution is asserted. Permanent source ID 30003378 and code OWR-15210-004 checked against the complete supplied registry and current repository. Distinct upstream records, including repeated questions and broad or damaged extracts, are retained for complete coverage of the source CS category. A repeated record is not a new independent conjecture. Complete mathematical-equivalence and current-literature review has not been performed. Preparation and final review follow contribution prompts C and D, with source gaps disclosed. Structural and rendering checks do not establish a new theorem or open status.
\label{end:30003378}
\end{document}
